
%\needspace{4\baselineskip}
\color{uniblue}{\subsection{Защита от снижения частоты (ЗСЧ)}}
\color{black}

\begin{enumerate}[label=\arabic{section}.\arabic{subsection}.\arabic*, labelsep=4pt, leftmargin=0pt, itemindent=57pt]

\item
Защита от снижения частоты предназначена для контроля снижения уровня частоты в сети. ЗСЧ может быть использована для выполнения отключения присоединений при недопустимо низком уровне частоты. В устройстве защита от снижения частоты представлена независимыми функциями (ступенями) «ЗСЧ 1» и «ЗСЧ 2».

\item
Основные характеристики функций:
\begin{itemize}
    \item погрешность ИО частоты не более 20 мГц;
    \item погрешность срабатывания ИО скорости изменения частоты не более 150 мГц;
    \item погрешность ИО частоты в рабочем диапазоне при напряжении входного сигнала 100 В не более 10 мГц; 
    \item время срабатывания и возврата ИО частоты не более 60 мс;
    \item время срабатывания ИО скорости снижения частоты не более 200 мс.
\end{itemize}

\item
ЗСЧ реагирует на снижение частоты ниже уставки. Срабатывание ИО функций формирует выходные сигналы «ЗСЧ 1: ИО F<», «ЗСЧ 2: ИО F<» для ЗСЧ 1, ЗСЧ 2 соответственно.

\item
Срабатывание функций происходит после отсчета выдержки времени, задаваемая уставкой Т1. По истечению выдержки времени таймера Т1, частота в течение этого времени постоянно оставалась ниже уставки срабатывания, а также отсутствовали сигналы блокировки, защита срабатывает и активируется соответствующие выходные сигналы функций «ЗСЧ 1: Срабатывание», «ЗСЧ 2: Срабатывание».

\item
Если до окончания времени таймера выдержки времени на срабатывание частота становилась выше уставки возврата органа частоты ЗСЧ, то активируется таймер задержки на возврат ИО Т2. Таймер задержки на возврат дает возможность таймеру выдержки времени на срабатывание продолжить работу, даже в случае, если частота кратковременно становилась выше уставки.

\item 
Блокировка функций осуществляется по следующим условиям:
\begin{itemize}
    \item снижение одного из линейных напряжений ниже уставки;
    \item снижение скорости изменения частоты ниже уставки, если уставка SGF2 введена;
    \item наличие входного дискретного сигнала внешней блокировки «Блокировка ЗСЧ».
\end{itemize}

\item
Ввод функций в работу на этапе параметрирования устройства осуществляется своим программным переключателем SGF1 «Ввод\_функции». Информация о введенной функции отображается выходным сигналом «ЗСЧ 1(2): Ввод».

\item
Функции могут быть выведены из работы оперативно путем активации сигнала «ОВ ЗСЧ 1(2)», что характеризуется активным состоянием сигнала «ЗСЧ 1(2): Оперативный вывод».

\item
При активном состоянии входного сигнала «ЗСЧ 1 на сигн», «ЗСЧ 2 на сигн» действие соответствующей функции ЗСЧ переводится на сигнал.

\item
Алгоритм работы функции «ЗСЧ 1» выполнен в соответствии с рисунком \ref{zsch:img1}. Алгоритм работы функции «ЗСЧ 2» аналогичный. В таблице \ref{zsch:tbl1} приведены параметры, необходимые для настройки функций.

\vspace{3mm}
\begin{figure}[H]
\centering
\includegraphics[width=1\textwidth,height=1\textheight,keepaspectratio]{\fbpath/116. ЗСЧ/_latex/img/zsch.pdf}
\captionof{figure}{Функциональная схема логики ЗСЧ 1}\label{zsch:img1}
\end{figure}

\small
\begin{longtable}{|>{\centering\arraybackslash}m{5.3cm}|>{\centering\arraybackslash}m{3.3cm}|>{\centering\arraybackslash}m{4.2cm}|>{\centering\arraybackslash}m{1.8cm}|>{\centering\arraybackslash}m{1cm}|}
\caption{Параметры для настройки ЗСЧ\hfill\vspace{-0.5\baselineskip}}\label{zsch:tbl1}\\ 
\hline
\rowcolor{gray!30}
Параметр (Параметр на ИЧМ) & Условное обозначение на схеме & Значение/ Диапазон & Единица измерения & Шаг \\ 
\hline
\endfirsthead
\caption*{\hspace{3pt}\emph{Продолжение таблицы \ref{zsch:tbl1}\hfill\vspace{-0.5\baselineskip}}} \\ % сделано по ГОСТ 2.105 п.6.8.7
\hline
\rowcolor{gray!30}
Параметр (Параметр на ИЧМ) & Условное обозначение на схеме & Значение/ Диапазон & Единица измерения & Шаг \\ 
%\hline
\endhead
%\hline
\endfoot
%\hline
\endlastfoot
%==+t1*PTUF1|LVTUF1
\centering Ввод функции в работу (Ввод\_функции) & \centering SGF1 & \centering 0 = Не предусмотрено\\1 = Предусмотрено & \centering -- & \centering \arraybackslash -- \\
\hline
\centering Частота срабатывания (Fср) & \centering F< & \centering 45,0 ... 51,0 & \centering Гц & \centering \arraybackslash 0,1 \\
\hline
\centering Разность частот срабатывания и возврата (dFвоз) & \centering -- & \centering 0,04 ... 1,00 & \centering Гц & \centering \arraybackslash 0,01 \\
\hline
\centering Напряжение срабатывания (Uср) & \centering U< & \centering 2,0 ... 100,0 & \centering В & \centering \arraybackslash 0,1 \\
\hline
\centering Величина срабатывания по скорости изменения частоты (dFdtср) & \centering dF</dt & \centering 0,1 ... 15,0 & \centering Гц/с & \centering \arraybackslash 0,1 \\
\hline
\centering Коэффициент возврата по скорости изменения частоты (Kвоз) & \centering -- & \centering 0,20 ... 0,99 & \centering -- & \centering \arraybackslash 0,01 \\
\hline
\centering Выдержка времени на возврат (Tвоз) & \centering T2 & \centering 0,00 ... 10,00 & \centering с & \centering \arraybackslash 0,01 \\
\hline
\centering Блокировка по dF/dt (Блок\_по\_dF/dt) & \centering SGF2 & \centering 0 = Не предусмотрено\\1 = Предусмотрено & \centering -- & \centering \arraybackslash -- \\
\hline
\centering Выдержка времени срабатывания (Tср) & \centering T1 & \centering 0,00 ... 100,00 & \centering с & \centering \arraybackslash 0,01 \\
\hline
%===t1
\end{longtable}
\normalsize

\end{enumerate}
